\section{Mid-Term Examination --- Set C}
\label{sec:midterm_set_c}

\ExamHeader{Mid-Term Examination}{C}{60 Minutes}{30 Marks}{Units 1, 2 \& 3 (Languages, Fundamentals \& Hardware)}

\noindent
\textbf{Instructions:} All 30 questions are compulsory. Each question carries 1 mark. Select the single best answer.

\vspace{3mm}

\mcqitem{1}{In machine language, instructions are represented using:}{Alphabets only}{Binary digits ($0$ and $1$)}{Hexadecimal numbers only}{ASCII text}

\mcqitem{2}{What is the primary function of an Assembler?}{To convert high-level languages into machine code}{To convert assembly language mnemonics into machine code}{To execute code line-by-line}{To edit documents}

\mcqitem{3}{Which of the following is a high-level programming language?}{Machine code}{Java}{Assembly code}{Microcode}

\mcqitem{4}{Which phase of compilation groups source characters into tokens such as keywords and operators?}{Lexical Analysis}{Syntax Analysis}{Semantic Analysis}{Optimization}

\mcqitem{5}{In C/C++, header file expansions (like \texttt{\#include <stdio.h>}) are handled by the:}{Compiler}{Preprocessor}{Linker}{Loader}

\mcqitem{6}{The software that combines multiple compiled object files with libraries into a single binary executable is the:}{Preprocessor}{Linker}{Loader}{Debugger}

\mcqitem{7}{What does the operating system's loader do?}{Translates code to assembly}{Loads the executable from disk into main memory to run}{Finds syntax errors}{Edits code}

\mcqitem{8}{An error where the program runs but calculates $2 + 2 = 5$ because of a wrong formula is a:}{Syntax error}{Logic error}{Linker error}{Compiler warning}

\mcqitem{9}{What happens during the \textit{Fetch} stage of the CPU instruction cycle?}{The ALU performs addition}{The instruction is retrieved from memory into a CPU register}{The instruction is decoded}{The output is printed}

\mcqitem{10}{A finite, unambiguous sequence of steps designed to solve a computing problem is an:}{Mnemonic}{Algorithm}{Operand}{Opcode}

\mcqitem{11}{Who designed the Analytical Engine in the 1830s containing early concepts of memory and processor?}{Blaise Pascal}{Charles Babbage}{Herman Hollerith}{Alan Turing}

\mcqitem{12}{The stored-program concept was introduced in 1945 by:}{John von Neumann}{Charles Babbage}{Linus Torvalds}{Steve Jobs}

\mcqitem{13}{First-generation computers (1940--1956) were characterized by the use of:}{Transistors}{Vacuum tubes}{Integrated circuits}{Microprocessors}

\mcqitem{14}{Which hardware component replaced vacuum tubes in second-generation computers (1956--1963)?}{Microprocessors}{Transistors}{Integrated circuits}{Magnetic drums}

\mcqitem{15}{Integrated circuits (ICs) on single silicon chips were introduced in the:}{First generation}{Second generation}{Third generation}{Fifth generation}

\mcqitem{16}{The fourth generation of computers (starting in 1971) was defined by the development of:}{Vacuum tubes}{Single-chip microprocessors}{Relays}{Gears}

\mcqitem{17}{Fifth-generation computing focuses on which modern development area?}{Vacuum tubes}{Artificial Intelligence and parallel computing}{Punched cards}{Magnetic core memory}

\mcqitem{18}{The computer characteristic GIGO stands for:}{Guaranteed Input, Guaranteed Output}{Garbage In, Garbage Out}{Global Instruction, General Output}{General Input, Good Output}

\mcqitem{19}{Which CPU unit directs instruction execution and coordinates data flow between memory and the ALU?}{Arithmetic Logic Unit}{Control Unit (CU)}{System Clock}{Hard drive}

\mcqitem{20}{Which system bus transmits physical memory addresses from the processor to memory chips?}{Data Bus}{Address Bus}{Control Bus}{Power Bus}

\mcqitem{21}{Which system bus transfers actual instructions and data words between the CPU and memory?}{Control Bus}{Data Bus}{Address Bus}{Serial Bus}

\mcqitem{22}{High-performance computers used for scientific simulations, weather forecasting, and defense research are:}{Microcontrollers}{Supercomputers}{Workstations}{Embedded systems}

\mcqitem{23}{Large, fault-tolerant computers used by banks and government agencies for high-volume transactions are:}{Personal computers}{Mainframe computers}{Laptops}{Microcontrollers}

\mcqitem{24}{A small computer built directly into an automobile engine or washing machine to perform a dedicated task is a/an:}{Supercomputer}{Embedded computer}{Mainframe}{Server}

\mcqitem{25}{A computer that processes continuous physical variables like voltage or temperature is a/an:}{Digital computer}{Analog computer}{Hybrid computer}{Mainframe}

\mcqitem{26}{Why are GPUs preferred over CPUs for training deep learning models?}{GPUs have higher clock speeds for single threads}{GPUs contain thousands of cores designed for parallel matrix operations}{GPUs eliminate the need for RAM}{GPUs are non-volatile}

\mcqitem{27}{The CPU register that stores the memory address of the next instruction to fetch is the:}{Instruction Register (IR)}{Program Counter (PC)}{Accumulator}{Stack Pointer}

\mcqitem{28}{The CPU register that holds the instruction currently being executed or decoded is the:}{Program Counter (PC)}{Instruction Register (IR)}{Memory Buffer Register}{Status Register}

\mcqitem{29}{Which level of CPU cache is typically the fastest and closest to the processor execution core?}{L3 Cache}{L2 Cache}{L1 Cache}{Virtual memory}

\mcqitem{30}{Which printer technology uses powdered toner and an electrostatic drum to produce rapid documents?}{Inkjet printer}{Laser printer}{Dot matrix printer}{Thermal printer}

\newpage
\subsection*{Answer Key \& Explanations --- Mid-Term (Set C)}

\begin{answerkeytable}{Mid-Term Set C --- Answer Key}
\begin{tabular}{*{10}{c}}
\toprule
\textbf{Q} & \textbf{Ans} & \textbf{Q} & \textbf{Ans} & \textbf{Q} & \textbf{Ans} & \textbf{Q} & \textbf{Ans} & \textbf{Q} & \textbf{Ans} \\
\midrule
1 & \textbf{B} & 7 & \textbf{B} & 13 & \textbf{B} & 19 & \textbf{B} & 25 & \textbf{B} \\
2 & \textbf{B} & 8 & \textbf{B} & 14 & \textbf{B} & 20 & \textbf{B} & 26 & \textbf{B} \\
3 & \textbf{B} & 9 & \textbf{B} & 15 & \textbf{C} & 21 & \textbf{B} & 27 & \textbf{B} \\
4 & \textbf{A} & 10 & \textbf{B} & 16 & \textbf{B} & 22 & \textbf{B} & 28 & \textbf{B} \\
5 & \textbf{B} & 11 & \textbf{B} & 17 & \textbf{B} & 23 & \textbf{B} & 29 & \textbf{C} \\
6 & \textbf{B} & 12 & \textbf{A} & 18 & \textbf{B} & 24 & \textbf{B} & 30 & \textbf{B} \\
\bottomrule
\end{tabular}
\end{answerkeytable}

\begin{expbox}[Technical Solutions --- Mid-Term Set C]
\begin{itemize}[leftmargin=5mm]
    \item \textbf{Q.2 (B):} Assemblers translate assembly mnemonics into machine code on a 1-to-1 basis.
    \item \textbf{Q.9 (B):} The Fetch stage reads the instruction from memory into the CPU using the Program Counter.
    \item \textbf{Q.26 (B):} GPUs are massively parallel, processing thousands of data elements simultaneously, making them ideal for neural networks.
    \item \textbf{Q.27 (B):} The PC (Program Counter) points to the next instruction in sequence.
    \item \textbf{Q.30 (B):} Laser printers fuse powdered toner with heat and pressure onto paper.
\end{itemize}
\end{expbox}

\newpage
